% Nová příprava: 2026-10-28; 1. ročník SPŠ, etapa 13–24.
% Rozpis: 1_rocnik/tematicke_plany/Rozpis_vyuky_MAT_1_rocnik.tex, hodina 19.
% Původ: vlastní formulace výkladu a vlastní úlohy; metodická inspirace uvedenými zdroji.
% Kontrola, obsahové porovnání zdrojů a časy: planovani/DOKONCENI_PRIPRAV_13_24.md.
% Učebnice: ucebnice/01_Zakladni_poznatky/Mnoziny/Množiny.tex
% oddíly Rovnost množin, Podmnožiny, Množinové operace a Vennovy diagramy; dle tématu hodiny.
% Realisticky.cz: externi_zdroje/realisticky/01_Zakladni_poznatky/03_Mnoziny/010306_Reseni_slovnich_uloh_pomoci_Vennovych_diagramu_I/010306_Lekce.pdf
% Realisticky.cz: externi_zdroje/realisticky/01_Zakladni_poznatky/03_Mnoziny/010307_Reseni_slovnich_uloh_pomoci_Vennovych_diagramu_II/010307_Lekce.pdf
\documentclass[11pt,a4paper]{article}
% Spolecny styl pro ucitelske pripravy MAT (Work/MAT)
% Verze 1.1, 2026-10-23
\NeedsTeXFormat{LaTeX2e}
\ProvidesPackage{priprava}[2026/10/23 v1.1 Sablona ucitelskych priprav MAT]

% Zakladni typografie a cestina
\RequirePackage[T1]{fontenc}
\RequirePackage[utf8]{inputenc}
\RequirePackage[czech]{babel}
\RequirePackage{lmodern}
\RequirePackage{microtype}
\RequirePackage{amsmath,amssymb}
\RequirePackage{xcolor}
\RequirePackage{enumitem}
\RequirePackage[most]{tcolorbox}
\RequirePackage{tabularx,booktabs,array}
\RequirePackage{fancyhdr}
\RequirePackage{needspace}
\RequirePackage{geometry}
\RequirePackage[hidelinks]{hyperref}

\geometry{a4paper,top=17mm,bottom=18mm,left=23mm,right=23mm}
\definecolor{notegray}{gray}{0.39}
\setlength{\parindent}{0pt}
\setlength{\parskip}{2.5pt plus 1pt minus 0.5pt}
\setlist[itemize]{topsep=2pt,itemsep=1pt,parsep=0pt,leftmargin=1.45em}
\setlist[enumerate]{topsep=2pt,itemsep=2pt,parsep=0pt,leftmargin=1.65em}
\setcounter{secnumdepth}{0}
\renewcommand{\arraystretch}{1.12}

% Zahlavi: vlevo rocnik | globalni tema, vpravo cislo: nazev hodiny
% Zapati zustava zcela prazdne.
\pagestyle{fancy}
\fancyhf{}
\renewcommand{\headrulewidth}{0.35pt}
\renewcommand{\footrulewidth}{0pt}
\setlength{\headheight}{14pt}
\newcommand{\nastavPripravu}[4]{%
  \fancyhead[L]{\small #1\enspace\textbar\enspace #2}%
  \fancyhead[R]{\small #3: #4}%
}

% Minimalisticke cernobile boxy. Kazda definice zvlast.
\tcbset{commonlesson/.style={enhanced,breakable=false,colback=white,colframe=black,
  colbacktitle=white,coltitle=black,boxrule=0.45pt,
  arc=0.8mm,outer arc=0.8mm,left=3mm,right=3mm,top=1.6mm,bottom=1.6mm,
  toptitle=1mm,bottomtitle=0.7mm,fonttitle=\bfseries,
  attach boxed title to top left={xshift=2.5mm,yshift=-1.15mm},
  boxed title style={colback=white,colframe=white,boxrule=0pt,left=1mm,right=1mm}}}
\newtcolorbox{zadaniBox}[1][]{commonlesson,title={Zadání},#1}
\newtcolorbox{prikladBox}[1][]{commonlesson,title={Příklad},#1}
\newtcolorbox{definiceBox}[1][]{commonlesson,title={Definice},#1}

% Spolecny box pro zadani a jeho reseni, bez mezery mezi castmi.
% Pouziti:
%   \begin{prikladReseniBox}
%     Zadani ulohy.
%     \tcblower
%     \begin{align*} ... \end{align*}
%   \end{prikladReseniBox}
% Spodni cast ma jemny sedy podklad a delici linku; lze v ni sazet
% vice radku matematiky. Obe casti maji jen jeden spolecny ramecek.
\tcbset{solutionlesson/.style={
  commonlesson,
  enhanced jigsaw,
  breakable,
  colbacklower=black!3,
  segmentation style={solid,draw=black!25,line width=0.35pt},
  fontlower=\small\color{notegray},
  before lower={\noindent\textbf{Řešení:}\enspace},
  leftlower=3mm,rightlower=3mm
}}
\newtcolorbox{prikladReseniBox}[1][]{solutionlesson,title={Příklad},#1}
\newtcolorbox{zadaniReseniBox}[1][]{solutionlesson,title={Zadání},#1}

% Pomucky pro autora pripravy: sedive poznamky; \reseni pro starsi pripravy
\newcommand{\poznamka}[1]{\par\begingroup\small\color{notegray}\itshape
  \textnormal{\textbf{Poznámka pro učitele.}}\enspace #1\par\endgroup}
\newcommand{\reseni}[1]{\par\begingroup\small\color{notegray}
  \textbf{Očekávané řešení:}\enspace #1\par\endgroup}
\newcommand{\blok}[1]{\needspace{5\baselineskip}\subsection*{#1}}
\newcommand{\vyrok}[1]{\uv{#1}}

\usepackage{tikz}
\usetikzlibrary{patterns,calc,angles,quotes}
\nastavPripravu{1. ročník}{Množiny}{19}{Slovní úlohy a počty}
\begin{document}
\blok{V oblasti je počet, nikoli další prvek}
Uvažujeme modelové skupiny bez jmen konkrétních žáků.
Rozlišuj seznam prvků množiny a číslo udávající jejich počet.
V diagramu můžeme místo jednotlivých prvků psát počty, musíme to však výslovně říci.
\begin{definiceBox}
Pro konečnou množinu $A$ značí \textbf{počet prvků} $|A|$ nezáporné celé číslo.
Například $|\{0,2,5\}|=3$ a $|\varnothing|=0$.
Svislé čáry zde označují počet, nikoli absolutní hodnotu čísla.
\end{definiceBox}

\begin{prikladReseniBox}[title={P1: dva kroužky},breakable=false]
Ve skupině 30 studentů navštěvuje robotiku 18 a programování 14 studentů.
Oba kroužky navštěvuje 6 studentů. Kolik chodí jen na robotiku,
jen na programování, alespoň na jeden a na žádný?
\tcblower
Označ $A$ robotiku a $B$ programování.
Jen $A$: $18-6=12$, jen $B$: $14-6=8$.
Alespoň jeden: $12+6+8=26$; žádný: $30-26=4$.
\begin{center}
\begin{tikzpicture}[scale=.72,color=black]
\draw (-2.5,-1.55) rectangle (2.65,1.7);
\draw (-.8,0) circle(1.25); \draw (.8,0) circle(1.25);
\node at (-1.2,1.4) {$A$}; \node at (1.2,1.4) {$B$};
\node at (2.25,1.4) {$U$};
\node at (-1.1,0) {$12$}; \node at (0,0) {$6$};
\node at (1.1,0) {$8$}; \node at (2.15,-1.2) {$4$};
\end{tikzpicture}
\end{center}
Čísla v obrázku jsou \emph{počty studentů v oblastech}.
Kontrola: $12+6=18$, $8+6=14$, $12+6+8+4=30$.
\end{prikladReseniBox}

Pro dvě konečné množiny platí
\[
|A\cup B|=|A|+|B|-|A\cap B|.
\]
Při součtu $|A|+|B|$ jsme společné prvky započítali dvakrát, proto jeden
jejich počet odečteme. Sčítat bez opravy lze počty disjunktních oblastí.

\begin{center}
\begin{tabularx}{\linewidth}{@{}l X@{}}
\toprule
Slovní údaj & Které oblasti zahrnuje \\ \midrule
Obě vlastnosti & Společnou část $A\cap B$.\\
Alespoň jedna & Jen $A$, společnou část i jen $B$.\\
Právě jedna & Jen $A$ nebo jen $B$, bez společné části.\\
Nejvýše jedna & Právě jednu, ale také žádnou.\\
Žádná & Oblast v $U$ mimo $A\cup B$.\\ \bottomrule
\end{tabularx}
\end{center}

\begin{prikladReseniBox}[title={P2: hledáme překryv},breakable=false]
Z 48 součástek má 17 poškrábaný povrch a 19 chybný rozměr.
Alespoň jednu z těchto vad má 29 součástek. Urči počet s oběma vadami,
jen s jednou a bez těchto vad.
\tcblower
Překryv $17+19-29=7$. Jen povrch: $17-7=10$, jen rozměr: $19-7=12$.
Právě jedna vada: $10+12=22$, bez obou uvedených vad: $48-29=19$.
Kontrola $10+7+12+19=48$. „Bez těchto vad“ neznamená důkaz, že součástka
nemůže mít jinou vadu, kterou zadání vůbec nesledovalo.
\end{prikladReseniBox}
\poznamka{Jádro P1, odvození vzorce, tabulka, P2 a S1 a), c), d):
asi 38 min práce, organizace a závěr 7 min. S1 b) je navazující zásoba.
Nejdříve kontroluj překlad textu do oblastí, teprve potom aritmetiku.
Soustavy několika rovnic nejsou nutným předpokladem této hodiny.}

\blok{Samostatné procvičení}
\begin{zadaniReseniBox}[title={S1: nejprve náčrt a význam údajů},breakable=false]
\begin{enumerate}[label=\alph*)]
\item Ve skupině 28 lidí ovládá program $A$ 16, program $B$ 11 a oba 5 lidí.
Urči oba výlučné počty, alespoň jeden, právě jeden a žádný.
\item Ze 40 výrobků má alespoň jednu ze dvou vlastností 31 výrobků.
Obě vlastnosti má 7. Pouze první má přesně dvojnásobný počet výrobků
než pouze druhou. Urči obě výlučné oblasti.
\item Mohou ve skupině 20 lidí platit počty $|A|=15$, $|B|=12$, $|A\cap B|=4$?
\item Kolik studentů z P1 navštěvuje nejvýše jeden kroužek?
\end{enumerate}
\tcblower
a) Jen $A$: $11$, jen $B$: $6$, alespoň jeden $22$, právě jeden $17$,
žádný $6$. Kontrola $11+5+6+6=28$.
b) Právě jednu vlastnost má $31-7=24$ výrobků.
Rozdělíme $24$ v poměru $2:1$, tedy pouze první $16$ a pouze druhou $8$.
Oba celkové počty jsou $23$ a $15$, mimo obě je $9$.
c) Ne, sjednocení by mělo $15+12-4=23$ lidí, více než celá skupina.
d) $30-6=24$, tedy $12+8+4$. Nejvýše jeden zahrnuje i studenty bez kroužku.
\end{zadaniReseniBox}

\blok{Volitelná rezerva: tři množiny}
\begin{zadaniReseniBox}[title={R1: postupuj od společného středu, 15–20 minut}]
Ve skupině 36 lidí platí $|A|=20$, $|B|=18$, $|C|=16$,
$|A\cap B|=9$, $|A\cap C|=8$, $|B\cap C|=7$ a $|A\cap B\cap C|=4$.
Párové průniky zahrnují i ty, kteří patří do všech tří.
Urči všech osm oblastí, počet v právě dvou množinách a počet mimo všechny.
\tcblower
Nejdříve společný střed: $4$.
Jen $A,B$: $9-4=5$, jen $A,C$: $8-4=4$, jen $B,C$: $7-4=3$.
Jen $A$: $20-5-4-4=7$; jen $B$: $18-5-3-4=6$;
jen $C$: $16-4-3-4=5$.
Ve sjednocení je $7+6+5+5+4+3+4=34$, mimo všechny $36-34=2$.
Právě ve dvou je $5+4+3=12$; alespoň ve dvou by bylo $12+4=16$.
Kontrola původních údajů např. $|A|=7+5+4+4=20$ a $|A\cap B|=5+4=9$.
\end{zadaniReseniBox}
\poznamka{Rezervu se třemi množinami zadávej až po zvládnutí dvou množin.
Při málu času ji přesuň na dobrovolné domácí rozšíření.}

\blok{Domácí procvičení}
\begin{zadaniReseniBox}[title={D1: 12–18 minut}]
\begin{enumerate}[label=\alph*)]
\item Z 32 lidí používá při cestě do školy vlak 20, autobus 17, obojí 9.
Urči pouze vlak, pouze autobus, alespoň jeden a žádný z těchto dopravních prostředků.
\item Z 50 dílů má vadu $A$ 21 dílů, vadu $B$ 18 dílů a ani jednu 19 dílů.
Kolik má obě vady a kolik právě jednu?
\item Kolik dílů z b) má nejvýše jednu z uvedených vad?
\item Proč nelze sčítat $|A|+|B|$, jsou-li některé prvky v obou množinách?
\end{enumerate}
\tcblower
a) Pouze vlak $11$, pouze autobus $8$, alespoň jeden $28$, žádný $4$.
b) Sjednocení $50-19=31$, obě $21+18-31=8$.
Pouze $A$: $13$, pouze $B$: $10$, právě jedna $23$.
c) $50-8=42=13+10+19$.
d) Společné prvky se započtou dvakrát; každý musí být ve sjednocení jen jednou.
\end{zadaniReseniBox}
\end{document}
